% =====================================================================
% backmatter/apendice-a.tex
% PT: Apêndice A. Os apêndices contêm material produzido pelo próprio
%     aluno que apoia o texto mas o tornaria demasiado longo (tabelas
%     completas, configurações, questionários, código). Depois de
%     \appendix em main.tex, cada \chapter passa a ser "Appendix A",
%     "Appendix B", etc. Para outro apêndice, copie este ficheiro
%     (ex.: apendice-b.tex) e junte o \input em main.tex.
% EN: Appendix A. Appendices hold material produced by the student that
%     supports the text but would make it too long (full tables,
%     configurations, questionnaires, code). After \appendix in main.tex,
%     each \chapter becomes "Appendix A", "Appendix B", etc. For another
%     appendix, copy this file (e.g. apendice-b.tex) and add its \input in
%     main.tex.
% =====================================================================
\chapter{Example Appendix}
\label{app:example}

% PT: Exemplo em inglês
% EN: English example
This appendix shows a table that uses the full text width with
\texttt{tabularx} and the \texttt{Y} column defined in
\texttt{config/commands.tex}. Refer to an appendix as any other chapter, for
example \cref{app:example}.

\begin{table}[htbp]
  \centering
  \caption{Example of a full-width table with tabularx.}
  \label{tab:appendix-example}
  \begin{tabularx}{\textwidth}{l Y}
    \toprule
    Parameter & Description \\
    \midrule
    Learning rate & Step size of the optimiser; see \cref{sec:method:setup}. \\
    Batch size    & Number of examples processed before each parameter update. \\
    Epochs        & Number of complete passes over the training data. \\
    \bottomrule
  \end{tabularx}
\end{table}

% PT: Exemplo em português
% EN: Portuguese example
\begin{otherlanguage}{portuguese}
Exemplo em português: a Tabela~\ref{tab:appendix-example} ocupa toda a
largura do texto. No corpo da dissertação, refira-se a este apêndice como
Apêndice~\ref{app:example}.
\end{otherlanguage}
